% BEGIN VOL08-EXPANSION VIII35
\section{Supplementary worked examples}
\label{sec:viii35-pedagogy}

\begin{example}[A degree-d map of the circle]\label{ex:viii35-ped-01}
For $f:S^1\to S^1$ of degree $d$,
\[
L(f)=\operatorname{tr}(f_*|H_0)-\operatorname{tr}(f_*|H_1)=1-d.
\]
If $d\ne1$, then $L(f)\ne0$, so the Lefschetz fixed-point theorem guarantees a fixed point.
\end{example}

\begin{example}[A degree-d map of a sphere]\label{ex:viii35-ped-02}
For $f:S^n\to S^n$ of degree $d$, only $H_0$ and $H_n$ contribute:
\[
L(f)=1+(-1)^n d.
\]
Thus the antipodal map has Lefschetz number $0$, consistent with its lack of fixed points.
\end{example}

\begin{example}[Linear torus maps]\label{ex:viii35-ped-03}
A torus map induced by an integer matrix $A\in M_n(\mathbb Z)$ acts on $H_1(T^n)$ by $A$ and on higher homology by exterior powers. The alternating trace identity gives
\[
L(f_A)=\det(I-A).
\]
A nonzero determinant therefore forces a fixed point.
\end{example}

\section{Graded supplementary exercises}
The following exercises progress from direct computation and construction to proof, hypothesis testing, application, and synthesis.

\subsection*{Standard computations and constructions}

\begin{exercise}[Definition]\label{exr:viii35-ped-01}
Define the Lefschetz number of a self-map $f:X\to X$ of a finite-type space.
\end{exercise}
\begin{hint}
Alternate traces on homology.
\end{hint}
\begin{solution}
$L(f)=\sum_i(-1)^i\operatorname{tr}(f_*:H_i(X;F)\to H_i(X;F))$ over a suitable field or finite free setting.
\end{solution}

\begin{exercise}[Identity map]\label{exr:viii35-ped-02}
What is $L(\operatorname{id}_X)$ for a finite CW complex?
\end{exercise}
\begin{hint}
The trace in each degree is the Betti number.
\end{hint}
\begin{solution}
$L(\operatorname{id}_X)=\chi(X)$.
\end{solution}

\begin{exercise}[Circle degree d]\label{exr:viii35-ped-03}
Compute $L(f)$ for a degree-$d$ map $S^1\to S^1$.
\end{exercise}
\begin{hint}
Use traces $1$ and $d$.
\end{hint}
\begin{solution}
$L(f)=1-d$.
\end{solution}

\begin{exercise}[Sphere degree d]\label{exr:viii35-ped-04}
Compute $L(f)$ for a degree-$d$ map $S^n\to S^n$.
\end{exercise}
\begin{hint}
Only degrees $0$ and $n$ contribute.
\end{hint}
\begin{solution}
$L(f)=1+(-1)^n d$.
\end{solution}

\begin{exercise}[Constant map]\label{exr:viii35-ped-05}
Compute the Lefschetz number of a constant self-map of a connected finite CW complex.
\end{exercise}
\begin{hint}
It is identity on $H_0$ and zero in positive degrees.
\end{hint}
\begin{solution}
$L(f)=1$.
\end{solution}

\subsection*{Proofs}

\begin{exercise}[Homotopy invariance]\label{exr:viii35-ped-06}
Prove homotopic self-maps have the same Lefschetz number.
\end{exercise}
\begin{hint}
Homotopic maps induce the same homology maps.
\end{hint}
\begin{solution}
The induced linear maps agree in every degree, so all traces and hence their alternating sum agree.
\end{solution}

\begin{exercise}[No-fixed-point contrapositive]\label{exr:viii35-ped-07}
What can be concluded if a map has no fixed points?
\end{exercise}
\begin{hint}
Use the contrapositive of the Lefschetz theorem.
\end{hint}
\begin{solution}
Under the theorem's hypotheses, it must have $L(f)=0$.
\end{solution}

\begin{exercise}[Torus determinant]\label{exr:viii35-ped-08}
Explain why $L(f_A)=\det(I-A)$ for a linear map of $T^n$.
\end{exercise}
\begin{hint}
Use traces on exterior powers.
\end{hint}
\begin{solution}
$H_k(T^n)\cong\Lambda^k\mathbb Z^n$, so the alternating sum $\sum(-1)^k\operatorname{tr}(\Lambda^kA)$ equals $\det(I-A)$.
\end{solution}

\begin{exercise}[Local index sum]\label{exr:viii35-ped-09}
State the relation between isolated fixed-point indices and the Lefschetz number.
\end{exercise}
\begin{hint}
Sum the local contributions.
\end{hint}
\begin{solution}
For isolated nondegenerate or suitably indexed fixed points, the sum of local fixed-point indices equals $L(f)$.
\end{solution}

\subsection*{Counterexamples and hypothesis tests}

\begin{exercise}[L=0 means no fixed point]\label{exr:viii35-ped-10}
Test: if $L(f)=0$, then $f$ has no fixed points.
\end{exercise}
\begin{hint}
The theorem gives only a sufficient condition when nonzero.
\end{hint}
\begin{solution}
False. Maps can have fixed points whose local indices cancel, or have zero Lefschetz number for other reasons.
\end{solution}

\begin{exercise}[Nonzero L counts fixed points exactly]\label{exr:viii35-ped-11}
Test: $|L(f)|$ always equals the number of fixed points.
\end{exercise}
\begin{hint}
Local indices may have signs or magnitudes.
\end{hint}
\begin{solution}
False. The Lefschetz number is an algebraic index sum, not generally a raw count.
\end{solution}

\begin{exercise}[L depends on chosen representative in homotopy class]\label{exr:viii35-ped-12}
Test: homotopic maps can have different Lefschetz numbers.
\end{exercise}
\begin{hint}
Induced homology maps are homotopy invariant.
\end{hint}
\begin{solution}
False. Lefschetz number depends only on the homotopy class.
\end{solution}

\subsection*{Applications and investigations}

\begin{exercise}[Fixed-point certificate]\label{exr:viii35-ped-13}
How can $L(f)$ provide a cheap certificate of a fixed point before solving $f(x)=x$ geometrically?
\end{exercise}
\begin{hint}
Compute finite-dimensional traces on homology.
\end{hint}
\begin{solution}
If the alternating trace is nonzero, the theorem guarantees at least one fixed point without locating it.
\end{solution}

\begin{exercise}[Dynamics on a torus]\label{exr:viii35-ped-14}
For a torus automorphism $A$, why is $\det(I-A)$ a useful dynamical diagnostic?
\end{exercise}
\begin{hint}
Nonzero determinant is the Lefschetz number.
\end{hint}
\begin{solution}
If $\det(I-A)\ne0$, a fixed point is forced; the determinant can be computed directly from the integer matrix defining the map.
\end{solution}

\subsection*{Challenge problems}

\begin{exercise}[Antipodal sphere]\label{exr:viii35-ped-15}
Check the Lefschetz number of the antipodal map on $S^n$.
\end{exercise}
\begin{hint}
Its degree is $(-1)^{n+1}$.
\end{hint}
\begin{solution}
$L=1+(-1)^n(-1)^{n+1}=0$, consistent with the fact that the antipodal map has no fixed points.
\end{solution}

\begin{exercise}[Lefschetz synthesis]\label{exr:viii35-ped-16}
Explain the chain of ideas from a geometric self-map to a fixed-point conclusion.
\end{exercise}
\begin{hint}
Pass to homology, take traces, then invoke the theorem.
\end{hint}
\begin{solution}
A self-map induces finite-dimensional endomorphisms on homology; their alternating trace is the homotopy-invariant Lefschetz number; if it is nonzero, the fixed-point theorem forces an intersection of the graph of $f$ with the diagonal, hence a fixed point.
\end{solution}

% END VOL08-EXPANSION VIII35
